\scanpage{19}
\uline{Proposition 2} Soit $G$ un schéma en groupes sur $S$.
Pour qu'il soit d'ordre infinitésimal $\nu$, i.e.
$U_\nu(G)=G$, il f. et s. que $G$ soit \uline{affine} sur $S$, et défini
par une Algèbre augmentée $A$
\deleted{[à \unsure{comultiplication} \unsure{diagonale donnée}]}
\deleted{\inserted{et tel que $x\in I\unsure{\Rightarrow}x^{p^\nu}=0$}}
dont l'\uline{I}déal d'augmentation \deleted{est de}
\inserted{a tous les éléments de} puissances $p^\nu$-ièmes nulles.
D'autre part, pour $G$ quelconque, $U_\nu(G)$ est
\deleted{can. isomorphe}\inserted{identique}
au voisinage infinitésimal d'ordre $p^\nu$ de la section unité.
\margin{N.B. Ne fait intervenir que la section marquée de $G$,
et pas la structure de groupe}
\mspage{18}
